\documentclass[11pt]{amsart}
\usepackage{amsmath, amssymb, amsthm}
\usepackage{geometry}
\usepackage{mathtools}
\geometry{margin=1in}
\usepackage[colorlinks=true, linkcolor=red, citecolor=red, urlcolor=blue]{hyperref}
\usepackage{perpage}
\MakePerPage{footnote}
\renewcommand{\thefootnote}{\fnsymbol{footnote}}

\theoremstyle{plain}
\newtheorem{theorem}{Theorem}[subsection]
\newtheorem{lemma}[theorem]{Lemma}
\newtheorem{proposition}[theorem]{Proposition}
\newtheorem{corollary}[theorem]{Corollary}

\theoremstyle{definition}
\newtheorem{definition}[theorem]{Definition}
\newtheorem{example}[theorem]{Example}
\newtheorem{exercise}[theorem]{Exercise}

\theoremstyle{remark}
\newtheorem{remark}[theorem]{Remark}

\numberwithin{equation}{section}

\begin{document}

\title[SDEs and Generative modeling]{A Treatise on Stochastic Differential Equations and their application to Generative modeling}

\author{Gowri Prakash}
\address{Department of Physics, Indian Institute of Science Education and Research, Tirupati, India.}
\email{gowriprakash.acad@gmail.com}

\thanks{This expository work was developed during my reading project at the Harish-Chandra Research Institute, under the supervision of Dr. Prasenjit Sen, in Summer 2026.}

\subjclass[2020]{}
\keywords{}

\begin{abstract}
This article explores the theoretical framework of stochastic differential equations (SDEs) and their application to generative modeling, particularly in the context of score-based generative models using SDEs. We delve into the construction of continuous-time SDEs, the dynamics of probability densities, and the reverse process that enables sampling from complex distributions. 
\end{abstract}

\maketitle

\section{Introduction}

A finite-dimensional It\^{o} stochastic differential equation (SDE) is of the form
\begin{equation}
    dX_t = \mu(t, X_t)\,dt + \sigma(t, X_t)\,dW_t.
\end{equation}
Now the differential form is strictly a shorthand notation for the integral form and the differential $dX_t$ is not a well-defined object in the classical sense.
The integral form is
\begin{equation}
    X_t = X_0 + \int_0^t \mu(s, X_s)\,ds + \int_0^t \sigma(s, X_s)\,dW_s.
    \label{eq:sde soln}
\end{equation}
Here $X_t$ is the state variable (which is a random variable; refer to Appendix~\ref{app:A}), $\mu:\mathbb{R}\times \mathbb{R}^n \to \mathbb{R}^n$ is the drift coefficient, and $\sigma:\mathbb{R}\times \mathbb{R}^n \to \mathbb{R}^{n \times m}$ is the diffusion coefficient when $\{W_t\}_{t \geq 0}$ is an $m$-dimensional Wiener process (refer to Appendix~\ref{app:B}) i.e., $W_t: \Omega \to \mathbb{R}^m$.\\

Now one may ask, the state variable $X_t$ here is not a general random variable whose codomain can be any measurable space, but is a very special random variable whose codomain is $\mathbb{R}^n$. You're right, SDEs are space specific. The SDE given, as mentioned is a finite dimensional It\^{o} SDE, and the state variable $X_t$ is a $\mathbb{R}^n$-valued random variable. However, there are also infinite-dimensional SDEs where the state variable can take values in function spaces or other infinite-dimensional spaces. The choice of $\mathbb{R}^n$ is common in many applications, but the theory of SDEs can be extended to more general settings as well.\\

Now $\int_0^t \mu(s, X_s)\,ds$ (say $I_t$) is a standard Lebesgue integral
\begin{equation}
    I_t(\omega)=\left( \int_0^t \mu(s, X_s)\,ds \right)(\omega) = \int_0^t \mu(s, X_s(\omega))\,ds := \int_{[0,t]} \mu(s, X_s(\omega)) \, d\lambda
    \label{dettol}
\end{equation}
for almost every $\omega \in \Omega$. Here, $\lambda$ is the Lebesgue measure on $[0, t]$. For a fixed $\omega$, the mapping $s \mapsto X_s(\omega)$ is a deterministic function of time, commonly referred to as a ``sample path''. Consequently, the mapping $s \mapsto \mu(s, X_s(\omega))$ defines a sample path of the composite stochastic process. By evaluating the Lebesgue integral of these sample paths individually for each $\omega$, we define the random variable $I_t : \Omega \to \mathbb{R}^n$ pointwise.\\

Equivalently, the Riemann construction of the same integral is 
\begin{equation}
    I_t=\int_0^t \mu(s, X_s)\,ds = \lim_{n \to \infty} \sum_{i=0}^{n-1} \mu(t_i, X_{t_i}) (t_{i+1} - t_i)
    \label{yellow}
\end{equation}
assuming the mapping $s \mapsto \mu(s, X_s)$ defines an almost surely (a.s.) continuous stochastic process.

\remark{Almost surely (a.s.) is the probability theoretic analogue of almost everywhere (a.e.) in measure theory; i.e., a property holds a.s. if the probability measure of the set of all points where the property does not hold is zero. Here specifically, a stochastic process $\{X_t\}_{t \geq 0}$ (refer to Appendix~\ref{app:B}) is continuous almost surely if $P(\{\omega \in \Omega \mid  t \mapsto X_t(\omega) \text{ is discontinuous}\}) = 0$ where $P$ is the probability measure.}\\

If we manage to find an analytical solution for $I_t$ in closed form (i.e., for any $t$, we know the function $I_t:\Omega \to \mathbb{R}^n$ and it can be computed), then we're done. But if we can't, the only way out is computing the integral numerically. And \eqref{dettol} is incompatible for such a venture, because the Lebesgue integral construction partitions the state space, requiring full \textit{a priori} knowledge of the unknown sample path $X_t(\omega)$. Therefore we are forced to use the Riemann formulation in \eqref{yellow}. Since it partitions the time domain instead, we can compute the path sequentially, moving forward step by step from the known initial condition $X_0$ \\

Now comes the It\^{o} integral $\int_0^t \sigma(s, X_s)\,dW_s$ (say $J_t$). We start with the It\^{o} sum
\begin{equation}
    J_t^{(n)} = \sum_{i=0}^{n-1} \sigma(t_i, X_{t_i})(W_{t_{i+1}} - W_{t_i}).
\end{equation}
And we define $J_t$ as follows:
\begin{equation} 
    J_t = \int_0^t \sigma(s, X_s)\,dW_s := \lim_{n \to \infty} J_t^{(n)}.
\end{equation}    
I'm avoiding the description of the above convergence as it moves away from the scope of this exposition. For a detailed description, refer to Appendix~\ref{app:C}.\\
\remark{In equation (1.5), the integrand $\sigma(s, X_s)$ is evaluated at the left endpoint of each subinterval, which is a key feature of the It\^{o} integral. This preserves the predictability of the process, ensuring that $\sigma(t_i, X_{t_i})$ is independent of the forward Brownian motion $W_{t_{i+1}} - W_{t_i}$. \\

And that's how we compute the evolution of the random variable as a function driven by an SDE. Because $X_t$ is defined on a probability space, this evolution theoretically dictates the exact evolution of its distribution via the pushforward measure. But say you're given the probability distribution of $X_0$ in $\mathbb{R}^n$ and asked about its time evolution driven by the SDE without giving the functional form of the state variable $X_0$, can you comment anything on that? Is there any equation governing such an evolution? That's a very specific question for me to make it seem like a natural question. So yes, such an equation exists. Okay, but how does all this have anything to do with generative modeling? Well, I haven't yet described what generative modeling even is. Let's get into the next section where you'll be getting answers to all these questions.\\

\noindent \textbf{A short note on the appendices.} To ensure self-containment of the exposition, I tried to give definitions of almost all the mathematical constructs introduced throughout the article. That being said, there are things in the appendices that are used nowhere in the article. So the appendix section can be seen as a dictionary of sorts for the exposition, as it definitely contains things that go outside the subject matter.

\section{Generative modeling}

Generative modeling is a branch of machine learning that focuses on learning the \textit{underlying probability distribution} of a training dataset, in order to generate new samples from it, in addition to the training data.\\

What is the underlying distribution? Is there a `the' underlying distribution for a given dataset? Or is it something that depends on the model, i.e., how different models extrapolate information from the same training set? The concept of an \textit{empirical distribution} explains what \textit{generation} is. Say our training dataset is $\{x_i\}_{i=1}^m \subset \mathbb{R}^n$. The corresponding empirical distribution $\mathbb{P}_e$ (which is a measure on the target space, which enables integration with respect to it) of this dataset would then be
\begin{equation}
\mathbb{P}_e = \frac{1}{m}\sum_{i=1}^m \delta_{x_i}
\label{emp}
\end{equation}
where $\delta_{x_i}$ is the Dirac measure defined by
\[
\delta_{x_i}(A) = \mathbf{1}_A(x_i) = \begin{cases} 1 & \text{if } x_i \in A \\ 0 & \text{otherwise} \end{cases}
\]
for any Borel measurable set $A \in \mathcal{B}(\mathbb{R}^n)$\footnote{$\mathcal{B}(\mathbb{R}^n)$ is the borel sigma algebra of $\mathbb{R}^n$}.\\

Now we have a measure $\mathbb{P}_e$ on the target space $(\mathbb{R}^n,\mathcal{B}(\mathbb{R}^n))$ of the random variable. Drawing a sample from $\mathbb{P}_e$ always returns a point in $\{x_i\}_{i=1}^m$. So this distribution in some sense is the \textit{memorization} of the training dataset, as any sample drawn from it will always be an element of the training dataset. A generated sample $x^*$ is one such that $x^* \notin \{x_i\}_{i=1}^m$. Therefore to call a model `generative', we should be able to sample such $x^*$, which is what would be enabled by the underlying distribution.\\

So clearly, the empirical distribution is fixed for a given training dataset irrespective of the specific generative model in question. But what about the underlying distribution? Is it unique for a given dataset? No, it's not. Different generative models define different distributions over $\mathbb{R}^n$ for the same training dataset $\{x_i\}_{i=1}^m$. Say the training set is a collection of images of dogs. There's no universal underlying distribution for `dog-ness' which'll make a sample drawn from it likely to be dog-like. Different models predict dog-ness differently.\\

In the following sections, we'll see how sampling from such an underlying distribution is made possible in score-based generative models using SDEs.


\subsection{Score-based generative modeling using SDEs}

We start with our training dataset $\{x_i\}_{i=1}^m \subset \mathbb{R}^n$. Now think that all the elements of this training dataset are drawn from some unknown underlying distribution, say $p_0$\footnote{Using the same notation as ~\cite{song2020} for ease of referencing for the reader.} ($p_t$ is the marginal distribution of the state variable $X_t$ which is an element of the solution stochastic process $\{X_t\}_{t \geq 0}$ of the SDE). It's worth noting that the current picture of a distribution (i.e., of the underlying distribution $p_t$) is different from the one I gave for the empirical distribution (which was a measure). These distributions are defined on different domains. The former on $\mathbb{R}^n$ and the latter on $\mathcal{B}(\mathbb{R}^n)$. The former is a probability density function (PDF) while the latter is a probability measure. The latter is strictly a more general construct. But we'd be exclusively working with the former in this exposition as the papers we refer to, work with PDFs.\\

So what we have in hand to start with is a known state variable $X_0 : \Omega \to \mathbb{R}^n$ with an unknown probability density function $p_0$. 
The function $X_0$ is on which we employ the forward SDE $dX_t = \mu(t, X_t)\,dt + \sigma(t, X_t)\,dW_t$ to get a continuous time sequence $\{X_t\}_{t \geq 0}$\footnote{This notation is a doppelganger of the stochastic process notation, as the notation for a stochastic process $\{X_t\}_{t \geq 0}$ unfortunately lacks an indication of the probability distributions associated with the random variables.} where $X_t = X_0+I_t+J_t$ as demonstrated in \eqref{eq:sde soln}.
But really, this does not have the full information of the solution stochastic process (refer to Appendix~\ref{app:B}). 
The random variable $X_0$ that we are dealing with here, is really not a random variable in the measure-theoretic probability theory sense (as given in Appendix~\ref{app:A}), as our sample space $\Omega$ is not a probability space by itself, for us to be able to push the measure forward (i.e., the pushforward measure) to the target measurable space of $X_0$. 
In fact that's the whole problem of generative modeling, finding $p_0$. But the random variable not being such doesn't stop us from employing the SDE on it. 
As for the SDE, the state variable $X_0$ is just a $\mathbb{R}^n$-valued function that it morphs in time to get $X_t$ for all time $t \geq 0$. 
So how is the solution of the SDE, if exists, a stochastic process? As for $\{X_t\}_{t \geq 0}$ to be a stochastic process, we not only need the marginal probability distribution $p_t$ of $X_t$ for all $t \geq 0$, but also the joint probability distribution of $(X_{t_1}, X_{t_2}, \ldots, X_{t_n})$ for any finite collection of time points $0 \leq t_1 < t_2 < \ldots < t_n$. 
This is where the Fokker--Planck equation (refer to Appendix~\ref{app:D}) gets in the picture, the solution of which gives us the marginal probability distribution $p_t$ of $X_t$ for all $t \geq 0$ (i.e., a continuous time sequence $\{p_t\}_{t \geq 0}$), of course, if we know $p_0$ to start with, as the deterministic evolution of marginal probabilities needs an initial condition to have the determinism.\\

The FPE correspondence of the forward SDE (let's call it forward FPE) $dX_t = \mu(t, X_t)\,dt + \sigma(t, X_t)\,dW_t$ (which is an $n$-dimensional euclidean SDE) is 
\begin{equation}
    \frac{\partial p_t}{\partial t} = -\sum_{i=1}^n \frac{\partial}{\partial x_i}[\mu_i(t, x)p_t] + \frac{1}{2}\sum_{i,j=1}^n \frac{\partial^2}{\partial x_i \partial x_j}[(\sigma\sigma^T)_{ij}(t, x)p_t].
    \label{eq:FPE}
\end{equation}
So how is the FPE of any use when we don't know $p_0$? Because we need $p_0$ to uniquely determine $p_t$ for all $t \geq 0$. Not quite! Yes, when $t \rightarrow \infty$. We can engineer the coefficients $\mu$ and $\sigma$ of the forward SDE such that, no matter what $p_0$ is, the solution of the FPE $p_t$ converges to a known distribution as $t \rightarrow \infty$. Sure, but how is that of any use in tracing back to $p_0$? 
This is where Anderson's time-reversal theorem \cite{anderson1982} 
enters. It asserts that, for any finite time point $t=T$ of the forward process $\{{X_t}\}_{t \geq 0}$, there exists a time-reversed process $\{\tilde{X_t}\}_{t \in [0,T]}$ where
$\tilde{X}_t = X_{T-t}$, that satisfies a ``reverse-time SDE" initialized at the 
`known' distribution $p_T$, whose marginals trace back through 
$\{p_t\}_{t \in [0,T]}$ i.e., $\tilde{p}_t=p_{T-t}$ where $\tilde{p}_t$ is the marginal PDF of $\tilde{X}_t$, recovering $p_0$ at time $T$, since $\tilde{p}_T=p_0$.\\

Two questions now. (1) What is the ``reverse time SDE"? (2) Let's even say we have it, how can we trace back to $p_0$? No one can initialize the reverse-time SDE at infinity, as that's the only time point (It's not a point) where we can potentially know the marginal by coefficients engineering.\\

The reverse FPE in the canonical form (refer to Appendix~\ref{app:D}) reads
\begin{equation}
    \frac{\partial \tilde{p_{t}}}{\partial t} = -\sum_{i=1}^n \frac{\partial}{\partial x_i}
\left[\left(-\mu_i(\tau, x) +\frac{1}{\tilde{p_{t}}}\sum_{j=1}^n \frac{\partial}{\partial x_j}
[(\sigma\sigma^\top)_{ij}(\tau, x)\tilde{p_{t}}]\right)\tilde{p_{t}}\right]
+ \frac{1}{2}\sum_{i,j=1}^n \frac{\partial^2}{\partial x_i \partial x_j}
[(\sigma\sigma^\top)_{ij}(\tau, x) \tilde{p_{t}}]
    \label{eq:reverse-FPE}
\end{equation}
where $\tau=T-t$ and $t \in [0,T]$.\\

Applying the product rule to the second term inside the parentheses yields

\begin{equation}
    \frac{1}{\tilde{p_{t}}}\sum_{j=1}^n \frac{\partial}{\partial x_j}
    [(\sigma\sigma^\top)_{ij}(\tau, x)\tilde{p_{t}}] = \underbrace{\sum_{j=1}^n \frac{\partial}{\partial x_j}
    [(\sigma\sigma^\top)_{ij}(\tau, x)]}_{=\nabla_x \cdot (\sigma\sigma^\top)(\tau, x)}+\underbrace{\sum_{j=1}^n
    (\sigma\sigma^\top)_{ij}(\tau, x)\frac{\partial \ln\tilde{p_{t}}}{\partial x_j}}_{=(\sigma\sigma^\top)(\tau, x)\nabla_x \ln \tilde{p_{t}}(x)}.
\end{equation}
\begin{equation}
    \bar{\mu}(\tau, \tilde{X}_t):= -\mu(\tau, \tilde{X}_t) + \nabla_x \cdot (\sigma\sigma^\top)(\tau, \tilde{X}_t) + (\sigma\sigma^\top)(\tau, \tilde{X}_t)\nabla_x \ln \tilde{p_{t}}(\tilde{X}_t)
    \label{eq:mu-bar}
\end{equation}

Under some regularity conditions like $p_t(x) > 0$ (the general condition true for all PDFs being $p_t(x) \geq 0$) for all $x \in \mathbb{R}^n$ and $t \in [0,T]$,
$p_t \in C^{1,2}$\footnote{$C^{1,2}([0,T] \times \mathbb{R}^n)$ denotes the space of functions 
that are once continuously differentiable in $t$ and twice continuously 
differentiable in $x$.}, and $\nabla_x \ln p_t \in L^2(p_t)$\footnote{$L^2(p_t)$ denotes the space of functions 
$f : \mathbb{R}^n \to \mathbb{R}^n$ satisfying $\int_{\mathbb{R}^n} 
\|f(x)\|^2 p_t(x)\,dx < \infty$.}, 
the reverse FPE \eqref{eq:reverse-FPE} corresponds to the reverse-time 
SDE \cite{anderson1982}
\begin{equation}
    d\tilde{X}_t = \left[-\mu(\tau, \tilde{X}_t) + 
    \nabla_x \cdot (\sigma\sigma^\top)(\tau, \tilde{X}_t)\,+
    (\sigma\sigma^\top)(\tau, \tilde{X}_t)\nabla_x \ln \tilde{p_{t}}(\tilde{X}_t)\right]dt 
    + \sigma(\tau, \tilde{X}_t)\,d{W}_t.
    \label{eq:reverse-SDE}
\end{equation}
Using \eqref{eq:mu-bar},
\begin{equation}
    d\tilde{X}_t = \bar{\mu}(\tau, \tilde{X}_t)\,dt + \sigma(\tau, \tilde{X}_t)\,d{W}_t
\end{equation}
initialized at $\tilde{X}_0 \sim p_T$, and $\{{W}_t\}_{t \geq 0}$ is an $m$-dimensional Wiener process. The marginal of $\tilde{X}_t$ is $p_{T-t}$, so 
$\tilde{X}_T \sim p_0$.\\

The reverse SDE \eqref{eq:reverse-SDE} is completely determined by the 
forward SDE coefficients $\mu$, $\sigma$ and the score function 
$\nabla_x \ln p_t: \mathbb{R}^n \to \mathbb{R}^n$ for all $t \in [0,T]$ (since $\tilde{p_t}=p_{T-t}$). Unlike the forward SDE, the reverse SDE can't function independently. As $p_t$ only exists for $t \in [0,T]$ given by the forward FPE, so does $\tilde{p_t}$. Since $\mu$ and $\sigma$ are chosen by design, 
the only unknown is the score function.\\

Okay, now we got a way to go back from $t=T$ to $t=0$ by perfectly reversing the forward stochastic process $\{X_t\}_{t \geq 0}$ to get $\{\tilde{X_t}\}_{t \in [0,T]}$. But all we can know is the marginal at $t=\infty$. What do we do now? We consider a large but finite $T$ such that $p_T \approx p_\infty$ where $p_\infty$ is the known distribution we engineered the coefficients $\mu$ and $\sigma$ to converge to.\\

Let's get into the specifics, of the paper by Song et al.~\cite{song2020}. They restrict the general finite dimensional It\^{o} SDE to a specific class where the diffusion coefficient is strictly a scalar function of time (i.e., a constant field over $\mathbb{R}^n$ for each $t$), yielding a forward SDE of the form:
\begin{equation}
    dX_t = f(t, X_t)\,dt + g(t)\,dW_t
\end{equation}

This is a linear SDE with respect to the state variable $X_t$ (a linear SDE is one where the drift and diffusion coefficients are linear functions of the state variable). Quoting directly from the paper, ``The SDE has a unique strong solution as long as the coefficients
are globally Lipschitz in both state and time" and ``When
$f(\cdot,t)$ is affine, the transition kernel is always a Gaussian distribution, where the mean and variance are
often known in closed-forms and can be obtained with standard techniques".\\

The paper introduces two different SDEs namely the Variance Exploding SDE (VE SDE) and the Variance Preserving SDE (VP SDE).\\
The VE SDE being
\begin{equation}
    dX_t= \sqrt{\frac{d[\sigma^2(t)]}{dt}}dW_t 
\end{equation}
and the VP SDE being
\begin{equation}
    dX_t=-\frac{1}{2}\beta(t)X_tdt+\sqrt{\beta(t)}dW_t.
\end{equation}\\

Clearly, both have drifts affine in $X_t$ and therefore Gaussian transition/perturbation kernels\footnote{Transition/Perturbation kernel $p_{st}(x(t) \mid x(s))$ is just the conditional probability distribution at a time $t$ given the realization of the random variable at any time $s\leq t$, i.e., $X_s=x_s$.} as follows:

$$p_{0t}(x(t) \mid x(0)) = \begin{cases} \mathcal{N}(x(t); x(0), [\sigma^2(t) - \sigma^2(0)]I) & \text{(VE SDE)} \\ \mathcal{N}\left(x(t); x(0)e^{-\frac{1}{2} \int_0^t \beta(s)\mathrm{d}s}, [1 - e^{-\int_0^t \beta(s)\mathrm{d}s}]I\right) & \text{(VP SDE)} \end{cases}$$\\

The VE SDE and VP SDE are continuous time variants of the discrete diffusion models SMLD and DDPM respectively. The $\sigma(t)$ and $\beta(t)$ are the continuous time versions of the discrete variants usually denoted by the same symbols in the literature of SMLD and DDPM respectively.\\

So do these SDEs make the marginal converge to a known distribution when $t \to \infty$? For the VE SDE, No! As the variance grows unbounded as $t \to \infty$, the marginals of the VE SDE diverge, failing to converge to a proper probability distribution. But for the VP SDE, the marginals do converge, to the standard Gaussian $\mathcal{N}(0,I)$.\\

For the VE SDE, $\sigma(t)$ is engineered as a strictly increasing function where $\sigma(0) \approx 0$, and $\sigma(t)$ grows unboundedly as $t \to \infty$. 
Because the VE SDE lacks a restorative drift term, this unbounded growth causes the variance of the marginal distribution to have an unbounded increase (So `exploding' is sort of a misnomer and not to be taken in the usual sense of a function exploding at a point). 
While the process lacks a true stationary limit, setting the terminal time $T$ such that $\sigma^2(T)$ is massively larger than the data scale allows the terminal marginal to be practically approximated by $p_T \approx \mathcal{N}(0, \sigma^2(T)I)$.\\

For the VP SDE, $\beta(t)$ is a strictly positive function. 
To ensure the total variance remains bounded and the distribution converges, the schedule must satisfy the condition that $\int_0^t \beta(s)ds$ diverges as $t \to \infty$.
This integral condition forces the affine mean coefficient $e^{-\frac{1}{2}\int_0^t \beta(s)ds}$ to asymptotically decay to zero. Consequently, the transition kernel of the VP SDE naturally and mathematically converges to the standard Gaussian stationary distribution $\mathcal{N}(0, I)$. 
Okay, the transition kernel is just the conditional distribution. We want the marginals to converge to $\mathcal{N}(0, I)$. Here is a brief sketch of how it's implied.

\begin{align*}
p_t(x) &= \int_{\mathbb{R}^n} p_{0t}(x \mid x_0) p_0(x_0) dx_0 \implies \lim_{t \to \infty} p_t(x) = \lim_{t \to \infty} \int_{\mathbb{R}^n} p_{0t}(x \mid x_0) p_0(x_0) dx_0 \\
&\xRightarrow{\text{DCT\footnotemark}} \lim_{t \to \infty} p_t(x) = \int_{\mathbb{R}^n} \left( \lim_{t \to \infty} p_{0t}(x \mid x_0) \right) p_0(x_0) dx_0
\end{align*}
\footnotetext{DCT-Dominated Convergence Theorem}
As established, for the VP SDE, $\lim_{t \to \infty} p_{0t}(x \mid x_0) = \mathcal{N}(x; 0, I)$. Therefore
$$ \lim_{t \to \infty} p_t(x) = \int_{\mathbb{R}^n} \mathcal{N}(x; 0, I) p_0(x_0) dx_0 =\mathcal{N}(x; 0, I) \int_{\mathbb{R}^n} p_0(x_0) dx_0. $$

Finally, since the initial data distribution $p_0(x_0)$ is a valid probability density, its integral over the entire space is exactly $1$. Therefore
$$ \lim_{t \to \infty} p_t(x) = \mathcal{N}(x; 0, I).$$

While the above limit guarantees convergence to a standard Gaussian, practical implementation requires halting the forward process at some finite terminal time $T$. 
Consequently, the true terminal marginal $p_T(x)$ is never exactly equal to $\mathcal{N}(x; 0, I)$. 
When initializing the reverse SDE however, we draw samples from the standard Gaussian prior relying on our approximation that $p_T(x)=\mathcal{N}(x; 0, I)$. 
So it introduces an initialization error. Therefore technically, we won't be reaching the exact $p_0$ at the end of the reverse process $\{\tilde{X_t}\}_{t \in [0,T]}$, if at all, one such exists. \\

Okay, now we have a prior to initialize the reverse SDE. But what about the score function $\nabla_x\ln p_t$ for all $t \in (0,T]$ which we need in order to run the reverse SDE? This is where the Gaussian transition kernel comes handy. 

$$p_t(x) = \int_{\mathbb{R}^n} p_{0t}(x \mid x_0) p_0(x_0) dx_0 \implies \nabla_x p_t(x) = \int_{\mathbb{R}^n} \nabla_x p_{0t}(x \mid x_0) p_0(x_0) dx_0$$
Multiplying both sides by $\frac{\nabla_x \ln p_{0t}(x \mid x_0)}{\nabla_x \ln p_{0t}(x \mid x_0)} \cdot \frac{1}{p_t(x)}$ gives\\
$$\frac{\nabla_x p_t(x)}{p_t(x)} = \int_{\mathbb{R}^n} \frac{\nabla_x p_{0t}(x \mid x_0)}{p_{0t}(x \mid x_0)} \cdot \frac{p_0(x_0)p_{0t}(x \mid x_0)}{p_t(x)} dx_0.$$
Therefore
$$\nabla_x \ln p_t(x) = \int_{\mathbb{R}^n} \nabla_x \ln p_{0t}(x \mid x_0) \cdot \frac{p_0(x_0)p_{0t}(x \mid x_0)}{p_t(x)} dx_0.$$
By Bayes' rule,
\begin{equation}
\frac{p_0(x_0)p_{0t}(x \mid x_0)}{p_t(x)} = p_{t0}(x_0 \mid x)\footnote{Just like the transition kernel $f(x,x_0)=p_{0t}(x \mid x_0)$, $p_{t0}(x_0 \mid x)$ is a multivariable function given by \eqref{bayes-rule}. The transition kernel $p_{0t}(x \mid x_0)$ is the conditional distribution of the state $x_t$ given a realization at $t=0$. $p_{t0}(x_0 \mid x)$ is the conditional distribution of the initial state $x_0$ given a realization at time $t$.}.
\label{bayes-rule}
\end{equation}
Therefore
\begin{equation}
\nabla_x \ln p_t(x) = \int_{\mathbb{R}^n} \nabla_x \ln p_{0t}(x \mid x_0) p_{t0}(x_0 \mid x) dx_0.
\label{hello}
\end{equation}

Since the integral in \eqref{hello} is taken over the domain of $x_0$ with respect to the probability density $p_{t0}(x_0 \mid x)$, it takes the exact form of a conditional expected value. We can therefore rewrite the identity using standard expectation notation:
\begin{equation}
\mathbb{E}_{x_0 \mid x} \left[ \nabla_x \ln p_{0t}(x \mid x_0) \right] := \int_{\mathbb{R}^n} \nabla_x \ln p_{0t}(x \mid x_0) p_{t0}(x_0 \mid x) dx_0.
\label{identity}
\end{equation}

Substituting this definition yields the fundamental identity relating the marginal score to the transition kernel score:
$$ \nabla_x \ln p_t(x) = \mathbb{E}_{x_0 \mid x} \left[ \nabla_x \ln p_{0t}(x \mid x_0) \right] $$

To utilize this identity in practice, we employ the framework of Score Matching. The primary goal is straightforward: we need a practical way to approximate the true marginal score $\nabla_x \ln p_t(x)$ so that it can be substituted directly into the reverse SDE. We achieve this by introducing a time-dependent neural network $s_\theta(x, t)$ designed to mimic this vector field. The most intuitive way to force the network to match the true score is to minimize the expected squared difference between them. \\



The Explicit Score Matching (ESM) loss at a specific time $t$ is defined as this expected squared error over the marginal distribution:
$$ \mathcal{L}_{\text{ESM}}(\theta, t) := \mathbb{E}_x \left[ \left\| s_\theta(x, t) - \nabla_x \ln p_t(x) \right\|_2^2 \right]\footnote{$\| \cdot \|_2$ denotes the standard Euclidean ($L^2$) norm.}$$

Directly minimizing this objective requires evaluating the true marginal score $\nabla_x \ln p_t(x)$, which is unknown for all $t \in [0, T)$ even after our assumption that $p_T(x) = \mathcal{N}(x;0,I)$. The identity \eqref{identity} provides the mathematical bridge to bypass this requirement. Expanding the Explicit Score Matching (ESM) loss yields
\begin{align*}
\mathcal{L}_{\text{ESM}}(\theta, t) &= \mathbb{E}_x \left[ \|s_\theta(x, t)\|_2^2 - 2\langle s_\theta(x, t), \nabla_x \ln p_t(x)\rangle + \|\nabla_x \ln p_t(x)\|_2^2 \right]\footnotemark.\\
&\implies \mathcal{L}_{\text{ESM}}(\theta, t) = \mathbb{E}_x \left[ \|s_\theta(x, t)\|_2^2 - 2\left\langle s_\theta(x, t), \mathbb{E}_{x_0 \mid x}[\nabla_x \ln p_{0t}(x \mid x_0)] \right\rangle \right] + C. 
\end{align*}
\footnotetext{$\langle \cdot, \cdot \rangle$ denotes the Euclidean inner product in $\mathbb{R}^n$.}
where $C = \mathbb{E}_x[\|\nabla_x \ln p_t(x)\|_2^2]=\int_{\mathbb{R}^n} \| \nabla_x \ln p_t(x) \|_2^2\cdot p_t(x) dx$.\\

Since $s_\theta(x, t)$ does not depend on $x_0$, it can be passed inside the inner expectation. Applying the law of total expectation ($\mathbb{E}_x \mathbb{E}_{x_0 \mid x}[\cdot] = \mathbb{E}_{x_0} \mathbb{E}_{x \mid x_0}[\cdot]$), we rewrite the objective as follows:
\begin{align*}
\mathcal{L}_{\text{ESM}}(\theta, t) &= \mathbb{E}_{x_0} \mathbb{E}_{x \mid x_0} \left[ \| s_\theta(x, t) \|_2^2 - 2 \langle s_\theta(x, t), \nabla_x \ln p_{0t}(x \mid x_0) \rangle \right] + C.\\
&\implies \mathcal{L}_{\text{ESM}}(\theta, t) = \mathbb{E}_{x_0} \mathbb{E}_{x \mid x_0} \left[ \| s_\theta(x, t) - \nabla_x \ln p_{0t}(x \mid x_0) \|_2^2 \right] + C' 
\end{align*}
where $C' = C - \mathbb{E}_{x_0} \mathbb{E}_{x \mid x_0} \left[ \| \nabla_x \ln p_{0t}(x \mid x_0) \|_2^2 \right]$.

$$ \mathcal{L}_{\text{DSM}}(\theta, t) := \mathbb{E}_{x_0} \mathbb{E}_{x \mid x_0} \left[ \left\| s_\theta(x, t) - \nabla_x \ln p_{0t}(x \mid x_0) \right\|_2^2 \right] $$

Because the losses $\mathcal{L}_{\text{ESM}}$ and $\mathcal{L}_{\text{DSM}}$ differ only by a constant independent of $\theta$, their gradients with respect to $\theta$ are identical. Consequently, the optimal parameters minimizing either objective are exactly the same:
$$ \arg\min_\theta \mathcal{L}_{\text{ESM}}(\theta, t) = \arg\min_\theta \mathcal{L}_{\text{DSM}}(\theta, t) $$

This establishes how we eliminate the need to evaluate the unknown marginal score, successfully replacing it with the computable score of the transition kernel, which we know in closed form.\\

To ensure the neural network learns to denoise across all noise scales, we extend this local objective by taking its expectation over the continuous time horizon $t \in (0, T]$. By weighting the loss at each time step with a strictly positive function $\lambda(t)$, we obtain the final continuous score matching objective as used in~\cite{song2020}:
\begin{equation} 
    \theta^* = \arg \min_\theta \mathbb{E}_t \left\{ \lambda(t) \mathbb{E}_{x_0} \mathbb{E}_{x \mid x_0} \left[ \left\| s_\theta(x, t) - \nabla_x \ln p_{0t}(x \mid x_0) \right\|_2^2 \right] \right\} 
    \label{score-obj}
\end{equation}
But still, evaluating this objective requires computing
$$ \mathbb{E}_{x_0}[f] = \int_{\mathbb{R}^n} f \, p_0(x_0) dx_0. $$
This computation strictly requires the unknown PDF $p_0(x)$. If you have not noticed, nowhere in our theoretical derivation up to this point have we actually used our training dataset $\{x_i\}_{i=1}^m$. This is exactly where the dataset enters the picture.\\

Because we cannot integrate over an unknown continuous density, we approximate the true expectation with the empirical expectation over our training set. By replacing the unknown measure with the empirical measure $\mathbb{P}_e$ defined in \eqref{emp}, the intractable integral reduces to a computable finite sum:
$$ \mathbb{E}_{x_0 \sim p_0}[f] \approx \mathbb{E}_{x_0 \sim \mathbb{P}_e}[f] = \frac{1}{m} \sum_{i=1}^m f(x_i). $$

By optimizing this empirical objective, we obtain optimal parameters $\theta^*$ such that $s_{\theta^*}(x, t) \approx \nabla_x \ln p_t(x)$ for almost all $x$ and $t$. With the intractable marginal score now successfully approximated by a neural network, we possess the exact vector field required to evaluate the reverse-time SDE.

In practice, sampling from this reverse process requires numerical discretization. Because the reverse drift term relies on a highly non-linear neural network $s_{\theta^*}(x, t)$, the reverse SDE cannot be integrated in closed form unlike the forward SDE, which admits an exact transition kernel that eliminates the need for sequential discretization. The standard approach employs numerical solvers, such as the Euler-Maruyama method, to partition the time horizon and iteratively simulate the stochastic trajectory backward. By initializing with pure noise $x_T \sim \mathcal{N}(0, I)$ and integrating back to $t \approx 0$, we transform the noise into a novel realization $x \sim p_0$, successfully generating new data from the underlying distribution.


\section{Conclusion}

Throughout this exposition, the initial data distribution $p_0(x)$ has been treated as a fixed, pre-existing continuous mathematical entity that the reverse stochastic process seeks to recover. However, the theoretical reality of generative modeling is far more nuanced. For real-world empirical datasets, a strictly continuous probability density $p_0(x)$ does not actually exist in nature; we are only ever provided a finite, discrete set of empirical observations, represented mathematically by the empirical measure $\mathbb{P}_e$.\\

The forward stochastic differential equation provides an analytical framework to map any abstract continuous distribution to stationary noise. Because we train the time-dependent neural network to approximate the reverse vector field using only the discrete empirical measure $\mathbb{P}_e$ (which means telling the score matching objective \eqref{score-obj} that $\mathbb{P}_e$ is, in fact, the initial distribution $p_0$), the model does not merely discover a hidden natural distribution. Rather, by optimizing the score matching objective over these discrete points, the neural network actively constructs a smooth, continuous probability density that plausibly interpolates the empty space between our training samples. In this sense, the generative model does not simply sample from a pre-existing $p_0$; it brings a continuous $p_0$ into existence.




\newpage
\appendix
\section{Probability Spaces and Random variables}
\label{app:A}

\begin{definition}[Measurable Space]
A measurable space is a pair $(X, \mathcal{F})$ where $X$ is a set and $\mathcal{F}$ is a $\sigma$-algebra of subsets of $X$.
\end{definition}

\begin{definition}[Measure Space]
A measure space $(X, \mathcal{F}, \mu)$ is a measurable space $(X, \mathcal{F})$ equipped with a measure $\mu:\mathcal{F} \to [0, \infty]$.
\end{definition}

\begin{definition}[Probability Space]
A probability space is a measure space $(\Omega, \mathcal{F}, P)$ where $\Omega$ is the sample space, $\mathcal{F}$ is a $\sigma$-algebra of subsets of $\Omega$ (called \textit{events}), and $P: \mathcal{F} \to [0,1]$ is a probability measure that is countably additive and $P(\Omega)=1$.
\end{definition}

\begin{definition}[Measurable Function]
A function $f: X \to Y$ between measurable spaces $(X, \mathcal{F})$ and $(Y, \mathcal{G})$ is measurable if for every $B \in \mathcal{G}$, the preimage $f^{-1}(B) \in \mathcal{F}$.
\end{definition}

\begin{definition}[$\sigma$-Algebra Generated by a Function]
$f: X \to (Y, \mathcal{G})$ is a function from a set $X$ to a measurable space $(Y, \mathcal{G})$. The $\sigma$-algebra generated by $f$ is $\sigma(f) = \{f^{-1}(B) : B \in \mathcal{G}\}$ (essentially making $f:(X,\sigma(f)) \to (Y,\mathcal{G})$ measurable).
\end{definition}

\begin{definition}[$\sigma$-Algebra Generated by a Collection of Sets]
$\mathcal{E}$ is a collection of subsets of a set $X$. The $\sigma$-algebra generated by $\mathcal{E}$, denoted $\sigma(\mathcal{E})$, is given by $\sigma(\mathcal{E}) = \bigcap \{\mathcal{F} \mid \mathcal{F} \text{ is a } \sigma\text{-algebra on } X \text{ and } \mathcal{E} \subset \mathcal{F}\}$.
\end{definition}

\begin{definition}[Pushforward Measure]
$f: (X, \mathcal{F}) \to (Y, \mathcal{G})$ is a measurable function and $\mu$ is a measure on $(X, \mathcal{F})$. The pushforward measure $f_* \mu$ on $(Y, \mathcal{G})$ is defined by $(f_* \mu)(B) = \mu(f^{-1}(B))$ for all $B \in \mathcal{G}$.
\end{definition}

\begin{definition}[Random Variable]
A random variable $X: (\Omega,\mathcal{F},P) \to (Y,\mathcal{G})$ is a measurable function defined on a probability space.
\end{definition}

\begin{definition}[Distribution of a Random Variable]
Let $X: \Omega \to Y$ be a random variable from a probability space $(\Omega, \mathcal{F}, P)$ to a measurable space $(Y, \mathcal{G})$. The distribution of $X$, denoted $P_X$, is the pushforward measure on $(Y, \mathcal{G})$ defined by $P_X(B) = P(X^{-1}(B))$ for all $B \in \mathcal{G}$.
\end{definition}

\begin{definition}[Joint Distribution]
$\{X_i\}_{i \in I}$ is a collection of random variables where $X_i : (\Omega, \mathcal{F}, P) \to (Y_i, \mathcal{G}_i)$. Their joint distribution is the distribution of the random variable $X = (X_i)_{i \in I}$, where $X: (\Omega, \mathcal{F}, P) \to (\prod_{i \in I} Y_i, \bigotimes_{i \in I} \mathcal{G}_i)$ is defined by $X(\omega) = (X_i(\omega))_{i \in I}$.
\end{definition}

\remark{The product space is $\prod_{i \in I} Y_i = \{ (y_i)_{i \in I} \mid y_i \in Y_i \}$. A coordinate projection map $\pi_j : \prod_{i \in I} Y_i \to (Y_j,\mathcal{G})$ is defined by $\pi_j((y_i)_{i \in I}) = y_j$. The product $\sigma$-algebra is $\bigotimes_{i \in I} \mathcal{G}_i = \sigma\left( \bigcup_{j \in I} \sigma(\pi_j) \right)$.}

\begin{definition}[Marginal Distribution]
Let $X = (X_i)_{i \in I}$ be a random variable constructed using the collection of random variables $\{X_i\}_{i \in I}$ with a given distribution (i.e., the joint distribution). For any $J \subset I$, the marginal distribution of the subcollection $X_J = (X_j)_{j \in J}$ is the distribution of $X_J: (\Omega, \mathcal{F}, P) \to (\prod_{j \in J} Y_j, \bigotimes_{j \in J} \mathcal{G}_j)$.
\end{definition}


\section{Stochastic Processes}
\label{app:B}
\begin{definition}[Stochastic Process]
    A stochastic process is a collection of random variables $\{X_t\}_{t \in T}$ defined on a common probability space $(\Omega, \mathcal{F}, P)$, where $T$ is an index set (often representing time).
\end{definition}

\begin{definition}[Filtration]
A filtration $\{\mathcal{F}_t\}_{t \geq 0}$ is an increasing family of $\sigma$-algebras such that $\mathcal{F}_s \subset \mathcal{F}_t$ for all $0 \leq s \leq t$.
\end{definition}

\begin{definition}[Natural Filtration]
    The natural filtration of a stochastic process $\{X_t\}_{t \in T}$ is the filtration $\{\mathcal{F}_t\}_{t \in T}$ defined by $\mathcal{F}_t = \sigma(X_s : s \leq t)$.
\end{definition}    

\begin{definition}[Markov Process]
    A stochastic process $\{X_t\}_{t \in T}$ is a Markov process if for any $s<t$ and any measurable set $A$, the following holds almost surely:
    \[
    P(X_{t} \in A \mid \mathcal{F}_s) = P(X_{t} \in A \mid \sigma(X_{s}))
    \]
    where $\mathcal{F}_s = \sigma(X_u : u \leq s)$.
\end{definition}

\remark{Notation. $P(X_{t} \in A \mid \mathcal{F}_s) := \mathbb{E}[1_A(X_t) \mid \mathcal{F}_s]$ where $1_A(X_t) : \Omega \to \{0, 1\}$ is the indicator function of the set $A$ defined as
 \[
 1_A(X_t)(\omega) = \begin{cases} 
  1 & \text{if } \omega \in A \\ 
  0 & \text{otherwise}.
\end{cases}
 \]}
Now this is a conditional expectation of a new random variable $1_A \circ X_t$ with respect to the $\sigma$-algebra $\mathcal{F}_s$. Refer Remark B.0.8.

\begin{definition}[Martingale]
    A stochastic process $\{X_t\}_{t \geq 0}$ is a martingale with respect to a filtration $\{\mathcal{F}_t\}_{t \geq 0}$ if for all $s < t$:
    \begin{enumerate}
        \item $X_s$ is $\mathcal{F}_s$-measurable.
        \item $\mathbb{E}[|X_t|] < \infty$.
        \item $\mathbb{E}[X_t \mid \mathcal{F}_s] = X_s$ almost surely.
    \end{enumerate}
\end{definition}  


\remark{Expectation of a random variable $X$ ($\mathbb{E}[X]$) is the Lebesgue integral of $X$ with respect to the probability measure $P$, i.e.\ $\mathbb{E}[X] = \int_\Omega X \, dP$. So checking the finiteness of $\mathbb{E}[|X_t|]$ is equivalent to checking the integrability of $X_t$. Very similar to the lebesgue integrable functions, where a function $f$ is lebesgue integrable if $\int |f| \, d\mu < \infty$, where $\mu$ is the Lebesgue measure.}

\remark{The conditional expectation $\mathbb{E}[X_t\mid\mathcal{F}_s]$ on the other hand, is the $P$-a.s.\ unique $\mathcal{F}_s$-measurable random variable $Z$ satisfying
\[
  \int_A Z \, dP = \int_A X_t \, dP \quad \forall A \in \mathcal{F}_s.
\]
Its existence follows from the Radon--Nikodym theorem.}


\begin{definition}[Wiener Process]
    A Wiener process (or Brownian motion) $\{W_t\}_{t \geq 0}$ is a stochastic process with the following properties:
    \begin{enumerate}
        \item $W_0 = 0$ almost surely.
        \item $W_t$ has independent increments: for $0 \leq s < t$, the increment $W_t - W_s$ is independent of $\mathcal{F}_s$, where $\mathcal{F}_s$ is the natural filtration of $\{W_t\}_{t \geq 0}$ 
        \item $W_t$ has stationary increments: for $0 \leq s < t$, the distribution of $W_t - W_s$ depends only on $t-s$.
        \item For each $t > 0$, $W_t$ is normally distributed with mean 0 and variance $t$: $W_t \sim N(0, t)$.
        \item The paths of $W_t$ are almost surely continuous.
    \end{enumerate}
\end{definition}  

\remark{The Wiener process is a Markov process and a martingale.}
\remark{An $\mathcal{F}_t$-Wiener process is one where an external filtration $\{\mathcal{F}_t\}_{t \geq 0}$ (to which the process is adapted) replaces the natural filtration in property (2).}

\begin{definition}[Predictable Process]
    A stochastic process $\{X_t\}_{t \geq 0}$ is predictable with respect to a filtration $\{\mathcal{F}_t\}_{t \geq 0}$ if for each $t \geq 0$, $X_t$ is $\mathcal{F}_{t-}$-measurable, where $\mathcal{F}_{t-} = \sigma(\mathcal{F}_s : s < t) = \sigma \left(\bigcup_{s < t} \mathcal{F}_s \right)$.
\end{definition}

\begin{definition}[Adapted Process]
    A stochastic process $\{X_t\}_{t \geq 0}$ is adapted to a filtration $\{\mathcal{F}_t\}_{t \geq 0}$ if for each $t \geq 0$, $X_t$ is $\mathcal{F}_t$-measurable, where $\mathcal{F}_t = \sigma(\mathcal{F}_s : s \leq t) = \sigma \left(\bigcup_{s \leq t} \mathcal{F}_s \right)$.
\end{definition}

\remark{A predictable process is always adapted, but an adapted process may not be predictable.}
\remark{A stochastic process is adapted to its natural filtration.}

\section{It\^{o} calculus}
\label{app:C}

While the framework of Itô calculus extends to integration with respect to general semimartingales, this exposition restricts the integrator strictly to the Wiener process.
Let $(\Omega, \mathcal{F}, \mathbb{P})$ be a probability space equipped with a filtration $\{\mathcal{F}_t\}_{t \ge 0}$, and let $\{W_t\}_{t \ge 0}$ be a standard one-dimensional $\mathcal{F}_t$-Wiener process. 

\subsection{Definition for Simple Processes}
We construct the integral for processes where evaluation is computationally exact. Let $\Delta = \{0 = t_0, t_1,\dots,t_N = T\}$ where $t_0 < t_1< \dots < t_N$ be a partition of the interval $[0, T]$. 

An adapted simple (step) process $C_t(\omega)$ is defined as:
\begin{equation}
    C_t(\omega) = \sum_{i=0}^{N-1} c_i(\omega) \mathbf{1}_{[t_i, t_{i+1})}(t)
\end{equation}
where each $c_i$ is $\mathcal{F}_{t_i}$-measurable and square-integrable ($\mathbb{E}[c_i^2] < \infty$). The $\mathcal{F}_{t_i}$-measurability guarantees that the value of the process over the interval $[t_i, t_{i+1})$ depends only on information available at or before time $t_i$.

The Itô integral of a simple process $C_t$ is defined by evaluating the integrand strictly at the left endpoint of each partition sub-interval:
\begin{equation}
    \int_0^T C_t \, dW_t := \sum_{i=0}^{N-1} c_i (W_{t_{i+1}} - W_{t_i})
\end{equation}

Because $c_i$ is determined at $t_i$ and the forward increment $(W_{t_{i+1}} - W_{t_i})$ is independent of $\mathcal{F}_{t_i}$ with expectation zero, the expectation of the integral is strictly zero.

\subsection{The Itô Isometry}
For the simple process $C_t$, the expectation of the squared integral maps directly to the standard Lebesgue integral of the squared process:
\begin{equation}
    \mathbb{E}\left[ \left( \int_0^T C_t \, dW_t \right)^2 \right] = \mathbb{E}\left[ \sum_{i=0}^{N-1} c_i^2 (t_{i+1} - t_i) \right] = \mathbb{E}\left[ \int_0^T C_t^2 \, dt \right]
\end{equation}
This identity, known as the Itô Isometry, provides the topological foundation to extend the integral to general processes via $L^2$-convergence.

\subsection{Extension to General Processes}
Let $\mathcal{L}^2_T$ denote the space of all real-valued, $\mathcal{F}_t$-adapted processes $f(t, \omega)$ that satisfy the integrability condition:
\begin{equation}
    \mathbb{E}\left[ \int_0^T f(t)^2 \, dt \right] < \infty
\end{equation}

For any general process $f \in \mathcal{L}^2_T$, there exists a sequence of simple processes $\{C^{(n)}\}_{n=1}^\infty$ that converges to $f$ in $L^2([0,T] \times \Omega)$:
\begin{equation}
    \lim_{n \to \infty} \mathbb{E}\left[ \int_0^T (f(t) - C^{(n)}_t)^2 \, dt \right] = 0
\end{equation}

By the Itô Isometry, the sequence of integrals $\int_0^T C^{(n)}_t \, dW_t$ forms a Cauchy sequence in $L^2(\Omega)$. The Itô integral of the general process $f(t)$ is defined as the unique mean-square limit of this sequence:
\begin{equation}
    \int_0^T f(t) \, dW_t := \lim_{n \to \infty} \int_0^T C^{(n)}_t \, dW_t \quad \text{(in } L^2\text{)}
\end{equation}

\subsection{Fundamental Properties}
For $f, g \in \mathcal{L}^2_T$ and constants $\alpha, \beta$:

\begin{enumerate}
    \item \textbf{Linearity:} 
    \begin{equation}
        \int_0^T (\alpha f(t) + \beta g(t)) \, dW_t = \alpha \int_0^T f(t) \, dW_t + \beta \int_0^T g(t) \, dW_t
    \end{equation}
    \item \textbf{Martingale Property:} The indefinite integral $M_t = \int_0^t f(s) \, dW_s$ is a continuous, square-integrable $\mathcal{F}_t$-martingale. Consequently:
    \begin{equation}
        \mathbb{E}\left[ \int_0^t f(s) \, dW_s \right] = 0
    \end{equation}
    \item \textbf{Quadratic Variation:} The quadratic variation of the stochastic process $M_t$ accumulates strictly as the time integral of the squared integrand:
    \begin{equation}
        [M, M]_t\footnote{The covariation $[X, Y]_t$ of two stochastic processes $\{X_t\}_{t \geq 0}$ and $\{Y_t\}_{t \geq 0}$ is $[X, Y]_t = \lim_{\max|t_i - t_{i-1}| \to 0} \sum_{i=1}^{n} (X_{t_i} - X_{t_{i-1}})(Y_{t_i} - Y_{t_{i-1}})$} = \int_0^t f(s)^2 \, ds
    \end{equation}
\end{enumerate}

\subsection{Itô's Lemma}
Because the quadratic variation of $W_t$ is non-zero ($[W,W]_t = t$), standard chain rule calculus fails. Second-order terms in the Taylor expansion do not vanish in the limit. 

Let $X_t$ be an Itô process governed by the stochastic differential equation:
\begin{equation}
    dX_t = \mu_t \, dt + \sigma_t \, dW_t
\end{equation}
where $\mu_t$ and $\sigma_t$ are adapted processes. Let $g(t, x)$ be a scalar function that is twice continuously differentiable in space and once in time, $g \in C^{1,2}([0, \infty) \times \mathbb{R})$. 

The differential of the transformed process $Y_t = g(t, X_t)$ is given by Itô's Lemma:
\begin{equation}
    dg(t, X_t) = \left( \frac{\partial g}{\partial t} + \mu_t \frac{\partial g}{\partial x} + \frac{1}{2} \sigma_t^2 \frac{\partial^2 g}{\partial x^2} \right) dt + \sigma_t \frac{\partial g}{\partial x} \, dW_t
\end{equation}


\section{Reverse Fokker-Planck equation}
\label{app:D}

The Fokker-Planck equation reads
\begin{equation}
    \frac{\partial p_t}{\partial t} = -\sum_{i=1}^n \frac{\partial}{\partial x_i} [\mu_i(t, x) p_t] + \frac{1}{2} \sum_{i,j=1}^n \frac{\partial^2}{\partial x_i \partial x_j} [(\sigma\sigma^\top)_{ij}(t, x) p_t].
    \label{sreejithaa}
\end{equation}

Define the reverse time parameter $\tau = T - t$. We define the marginal density of the reverse-time process $\tilde{X}_t$ as $\tilde{p}_t(x) = p_{T-t}(x) = p_\tau(x)$.

First, we differentiate $\tilde{p}_t$ with respect to the forward integration time $t$:
\begin{equation}
    \frac{\partial \tilde{p}_t}{\partial t} = \frac{\partial {p}_{T-t}}{\partial t} = \frac{\partial p_{\tau}}{\partial \tau}\cdot \frac{\partial \tau}{\partial t}   = -\frac{\partial p_\tau}{\partial \tau}
    \label{hell}
\end{equation}
\eqref{sreejithaa} is an identity holding for every time argument $s \in [0,T]$. Therefore
\[
\frac{\partial p_s(x)}{\partial s} = -\sum_{i=1}^n \frac{\partial}{\partial x_i}[\mu_i(s,x)p_s(x)]
+ \frac{1}{2}\sum_{i,j=1}^n \frac{\partial^2}{\partial x_i \partial x_j}[(\sigma\sigma^T)_{ij}(s,x)p_s(x)].
\]
As $\tau \in [0,T]$, we can comfortably substitute $s = \tau$, giving
\begin{equation}
\frac{\partial p_\tau}{\partial \tau} = -\sum_{i=1}^n \frac{\partial}{\partial x_i}[\mu_i(\tau,x)p_\tau]
+ \frac{1}{2}\sum_{i,j=1}^n \frac{\partial^2}{\partial x_i \partial x_j}[(\sigma\sigma^T)_{ij}(\tau,x)p_\tau(x)].
\label{hellobro}
\end{equation}
Now substituting \eqref{hell} in \eqref{hellobro} yields
\begin{equation}
    \boxed{\frac{\partial \tilde{p}_t}{\partial t} = \sum_{i=1}^n \frac{\partial}{\partial x_i} [\mu_i(\tau, x) \tilde{p}_t] - \frac{1}{2} \sum_{i,j=1}^n \frac{\partial^2}{\partial x_i \partial x_j} [(\sigma\sigma^\top)_{ij}(\tau, x) \tilde{p}_t]}
\end{equation}
The above equation is the reverse FPE.\\

To cast this equation into the canonical forward FPE form (which requires a negative drift divergence and a positive diffusion double-divergence), we add $\sum_{i,j=1}^n \frac{\partial^2}{\partial x_i \partial x_j} [(\sigma\sigma^\top)_{ij}(\tau, x) \tilde{p}_t]$ on both sides:
\begin{equation}
\begin{split}
    \frac{\partial \tilde{p}_t}{\partial t} &= \sum_{i=1}^n \frac{\partial}{\partial x_i} [\mu_i(\tau, x) \tilde{p}_t] - \sum_{i,j=1}^n \frac{\partial^2}{\partial x_i \partial x_j} [(\sigma\sigma^\top)_{ij}(\tau, x) \tilde{p}_t] \\
    &\quad + \frac{1}{2} \sum_{i,j=1}^n \frac{\partial^2}{\partial x_i \partial x_j} [(\sigma\sigma^\top)_{ij}(\tau, x) \tilde{p}_t]
\end{split}
\end{equation}
We group the first two terms by factoring out $-\sum_{i=1}^n \frac{\partial}{\partial x_i}$:
\begin{equation}
    \frac{\partial \tilde{p}_t}{\partial t} = -\sum_{i=1}^n \frac{\partial}{\partial x_i} \left[ -\mu_i(\tau, x) \tilde{p}_t + \sum_{j=1}^n \frac{\partial}{\partial x_j} [(\sigma\sigma^\top)_{ij}(\tau, x) \tilde{p}_t] \right] + \frac{1}{2} \sum_{i,j=1}^n \frac{\partial^2}{\partial x_i \partial x_j} [(\sigma\sigma^\top)_{ij}(\tau, x) \tilde{p}_t]
\end{equation}
Finally, to explicitly isolate the effective reverse drift from the density $\tilde{p}_t$, we multiply and divide the inner summation of the drift term by $\tilde{p}_t$. This recovers the final form:
\begin{equation}
\begin{split}
    \frac{\partial \tilde{p}_t}{\partial t} &= -\sum_{i=1}^n \frac{\partial}{\partial x_i} \left[ \left( -\mu_i(\tau, x) + \frac{1}{\tilde{p}_t} \sum_{j=1}^n \frac{\partial}{\partial x_j} [(\sigma\sigma^\top)_{ij}(\tau, x) \tilde{p}_t] \right) \tilde{p}_t \right] \\
    &\quad + \frac{1}{2} \sum_{i,j=1}^n \frac{\partial^2}{\partial x_i \partial x_j} [(\sigma\sigma^\top)_{ij}(\tau, x) \tilde{p}_t]
\end{split}
\end{equation}



\begin{thebibliography}{99}

\bibitem{song2020}
Y.~Song, J.~Sohl-Dickstein, D.~P. Kingma, A.~Kumar, S.~Ermon, and B.~Poole, 
\emph{Score-based generative modeling through stochastic differential equations}, 
arXiv preprint arXiv:2011.13456 (2020).
\bibitem{song2020}
Y.~Song and S.~Ermon,
\emph{Generative modeling by estimating gradients of the data distribution},
arXiv preprint arXiv:1907.05600 (2020).
\bibitem{anderson1982}
B.~D.~O.~Anderson,
\emph{Reverse-time diffusion equation models},
Stochastic Processes and their Applications, Vol.~12, No.~3, pp.~313--326, 1982.
\bibitem{luo2022}
C.~Luo,
\emph{Understanding diffusion models: a unified perspective},
arXiv preprint arXiv:2208.11970 (2022).
\bibitem{tao2011}
T.~Tao, \emph{An Introduction to Measure Theory}, Graduate Studies in Mathematics, Vol. 126, American Mathematical Society, Providence, RI, 2011.





\end{thebibliography}

\end{document}

